\chapter{Introduction}

% =====================================================================
% Se redacta EN ÚLTIMO LUGAR. Este capítulo fija el contrato de la tesis:
% las preguntas de investigación que aquí se enuncian son exactamente las
% que el Capítulo 7 debe responder una por una.
% =====================================================================

\section{Motivation}
\emph{Situar el problema: la telemetría vehicular continua exige entregar datos
de forma fiable y suficientemente fresca, pero el acceso celular permanente tiene
un coste energético que condiciona la autonomía del nodo embarcado. Introducir la
tensión central de la tesis entre disponibilidad, frescura del dato y energía.}

\section{Problem Statement}
\emph{Formular el problema con precisión: un vehículo con doble interfaz debe
decidir, en cada instante, si transmite ahora por la red celular o retiene el dato
esperando una oportunidad de conectividad inalámbrica futura, bajo capacidad de
memoria finita y plazos de validez del dato. Enunciar por qué esa decisión no es
trivial y por qué requiere predicción.}

\section{Objectives and Research Questions}
\emph{Objetivo general y objetivos específicos, seguidos de las preguntas de
investigación numeradas. Formularlas de modo que cada una sea contestable con
evidencia medida en el Capítulo 6; evitar preguntas que solo se puedan responder
cualitativamente.}

\section{Contributions}
\emph{Enunciar las contribuciones de forma explícita y numerada: el nodo
vehicular híbrido, el modelo energético dirigido por eventos, la capa de
predicción de oportunidades de conectividad, la política de descarga predictiva y
su extensión sensible a plazos, y la evaluación comparativa frente a las
referencias de tecnología única.}

\section{Structure of the Thesis}
\emph{Recorrido breve por los capítulos, indicando la función de cada uno dentro
del argumento y no solo su título.}
